\paragraph{My contribution.} I led the decision and optimisation part of the system, which turns risk scores and rules into a weekly contact plan. I built the case-based reasoner that estimates how much each offer would help each member, using 2,400 past cases, and shows the most similar past members as evidence for each recommendation. I built the genetic algorithm planner with its repair operator and greedy seeding, an exact CP-SAT planner that also serves as its benchmark, the baseline policies, and a policy evaluator that grades any plan against the scenario's ground truth. I ran experiments E3 to E5. On the app side I built the plan workflow: the plan builder, overrides checked live against the rules, approval, automatic assignment to staff with call sheets, scenario comparison, the staff contact list with outcome logging that feeds back into the case base, and the settings and health page. I also wrote the architecture, objectives, optimisation and E3--E5 parts of this report.

\paragraph{What I learned.} I came to this project from software engineering, new to AI, and that was challenging but exciting. My engineering habits helped: I prepared tests first and never reported results I had not tested. Several of the tests were checked by breaking the rule they guard on purpose, to prove they could fail.

The biggest surprise was the null test. My first version of the case-based reasoner scored well against the true effects. But when I shuffled which members received which offer, so that there was nothing left to find, it still ``found'' effects: it was measuring churn risk, not the offer. That taught me not to trust a good result until I am sure what it is measuring and that it behaves as expected when there is nothing to find.

Building CBR also changed how I understood it. I had thought it was the same as kNN. kNN only looks up the most similar cases. CBR works from a case base of previous cases, both successes and failures, and goes through a cycle: it reuses them, revises the estimate, and retains each new outcome for future use.

The experiments taught me that a better optimiser cannot make up for weak estimates. With perfect knowledge of the effects, the genetic algorithm kept as many members as a simple greedy planner or more. On the case-based estimates it kept fewer (16.7 against 21.0 at the default week), because a stronger search also finds the errors in its inputs. Retaining new cases has a trap too. The planner chooses whom to contact, so its own outcomes are not a fair sample, and only a random control group would let the system keep learning without bias.

\paragraph{Applying it.} I already work in product development supporting startups, and with the current boom in AI many of them are building AI systems. This project lets me understand how such systems work underneath, and how I can contribute beyond the engineering. I will ask what a result is actually measuring, test that a method finds nothing when there is nothing to find, and check the quality of the estimates before trusting an optimiser built on them.
